\LessonHeader{30}{One Pen, One Lamp, One Night}{一支笔，一盏灯，一个夜晚}{幽默随笔 · Humorous Essay}{话题\ 寒假作业·亲子\enspace|\enspace 新词\ \textbf{17}\enspace|\enspace 新义\ \textbf{22}\enspace|\enspace 短语\ \textbf{30}\enspace|\enspace 全文\ \textbf{898}\ 词}{Lesson 30\enspace·\enspace One Pen, One Lamp, One Night}
\begin{lessonbody}
\PARA{1}“What is \RV{}{Chinese}{} \RV{}{speed}{?} I am \PH{\NS{}{reporting}{}{新词性}{report}{/rɪˈpɔːt/}{vi.}{报道}\ live}{report live}{现场报道} from a \RV{}{winter}{} \RV{}{homework}{} \NS{}{rush}{.”}{新义}{}{/rʌʃ/}{n.[sing.]}{赶工；匆忙} “The \RV{}{holiday}{} is \RV{}{almost}{} \PH{\RVB{}{used}{}{use}\ up.}{use up}{用完} \SMK{y}With one \RV{}{pen}{} and one \RV{}{lamp}{,} the \RV{}{young}{} \RVB{}{flowers}{}{flower} of our \RV{}{country}{} can \PH{work\ \RVB{}{wonders}{}{wonder}}{work wonders}{创造奇迹} \RV{}{in}{} just one night...”\EMK{y}{1}
\YI{1}{一支笔，一盏灯，祖国的花朵们只用一夜就能创造奇迹……}

\PARA{2}For the last \RV{}{few}{} days, \RVB{}{videos}{}{video} like these have been all over my \RV{}{phone}{:} \RVB{}{kids}{}{kid} \NS{}{racing}{}{新义}{race}{/reɪs/}{vi.}{赶紧（做）} to \RV{}{finish}{} their \AG{}{winter}{} \AG{}{homework}{} before school starts. \RVB{}{Parents}{}{parent} \NW{}{complain}{,}{/kəmˈpleɪn/}{vi.}{抱怨} then \PH{sit\ up}{sit up}{熬夜} with them \RV{}{anyway}{.} In some homes the \RV{}{whole}{} \RV{}{family}{} \RVB{}{joins}{}{join} in, \NS{}{treating}{}{新义}{treat}{/triːt/}{vt.}{对待；看待} the \PH{\RV{}{half}{-}\AGB{}{finished}{}{finish}}{half-finished}{只做了一半的} \AG{}{homework}{} like an \RV{}{emergency}{.}

\PARA{3}At first I \RVB{}{laughed}{}{laugh} so \NS{}{hard}{.}{新义}{}{/hɑːd/}{adv.}{厉害地} Then my \RV{}{brain}{} \RV{}{suddenly}{} \RVB{}{woke}{}{wake} up: I \NS{}{pictured}{}{新词性}{picture}{/ˈpɪktʃə(r)/}{vt.}{想象} my \RV{}{own}{} child in those \AG{}{videos}{,} and I stopped \AGB{}{laughing}{}{laugh} at \RV{}{once}{.} So I \RVB{}{sent}{}{send} them to a good friend with two \RVB{}{sons}{.}{son} After all, two people \RVB{}{suffering}{}{suffer} is better \RV{}{than}{} one.

\PARA{4}\RV{}{Across}{} the \AG{}{country}{,} \RV{}{primary}{} school \RVB{}{students}{}{student} (and their \AG{}{parents}{)} \RV{}{seem}{} to \RV{}{follow}{} the \RV{}{same}{} \RV{}{pattern}{.} On the first day of the \AG{}{holiday}{,} they \RV{}{fly}{} \RV{}{free}{} like \RVB{}{birds}{.}{bird} \RV{}{Right}{} before school, they are \PH{\RVB{}{fighting}{}{fight}\ for\ their\ \RVB{}{lives}{.}{life}}{fight for one's life}{拼命挣扎}

\PARA{5}These \AG{}{kids}{} \RV{}{may}{} work \RVB{}{harder}{}{hard} \AG{}{than}{} \RV{}{Olympic}{} \RVB{}{champions}{.}{champion} \RV{}{Late}{} at night, when the \NW{}{skiing}{}{/ˈskiːɪŋ/}{n.[U]}{滑雪} \RV{}{star}{} \PN{}{Eileen Gu}{}{/ˈaɪliːn ˈɡuː/}{n.}{谷爱凌} is \RV{}{already}{} \RVB{}{enjoying}{}{enjoy} her ten \RVB{}{hours}{}{hour} of \RV{}{sleep}{,} they are \RV{}{still}{} at their \NWB{}{desks}{,}{desk}{/desk/}{n.[C]}{书桌} \PH{\AG{}{racing}{}\ \RV{}{against}{}\ time.}{race against time}{争分夺秒} \RV{}{During}{} the day, \RV{}{while}{} the \PH{\RV{}{ice}{}\ \NW{}{skating}{}{/ˈskeɪtɪŋ/}{n.[U]}{滑冰}}{ice skating}{滑冰；溜冰} \AG{}{champion}{} \PN{}{Wu Dajing}{}{}{n.}{武大靖} \NWB{}{skates}{}{skate}{/skeɪt/}{vi.}{滑冰} \AG{}{hard}{} on the \RV{}{ice}{,} they \RV{}{wish}{} their \RVB{}{hands}{}{hand} \RV{}{could}{} \RV{}{move}{} at a \AG{}{champion}{’s} \AG{}{speed}{.}

\PARA{6}The \AG{}{homework}{} can be done \RV{}{anywhere}{.} They write on the \NS{}{underground}{,}{新词性}{}{/ˈʌndəɡraʊnd/}{n.[sing.]}{地铁} in \RV{}{shopping}{} \RVB{}{centres}{,}{centre} in \RV{}{fast}{} food \RVB{}{restaurants}{}{restaurant} and on \RVB{}{planes}{.}{plane} No time? \PH{Make\ time!}{make time}{挤出时间} \RV{}{Fall}{} \RV{}{asleep}{?} Write \RV{}{with}{} your \RVB{}{eyes}{}{eye} \RVB{}{closed}{!}{close} \SMK{y}How much \AG{}{kids}{} and \AG{}{parents}{} can do \RVB{}{depends}{}{depend} only on how many \AG{}{hours}{} are \RVB{}{left}{}{leave} before school.\EMK{y}{2}
\YI{2}{孩子和家长能做多少，只取决于开学前还剩几个小时。}

\PARA{7}When a \AG{}{parent}{} \NW{}{carefully}{}{/ˈkeəfəli/}{adv.}{小心翼翼地} \RVB{}{points}{}{point} out, “I can’t \RV{}{even}{} read your \NW{}{writing}{,”}{/ˈraɪtɪŋ/}{n.[U]}{字迹} the child \AG{}{may}{} \PH{\RV{}{answer}{}\ \RV{}{right}{}\ back:}{answer back}{顶嘴；回嘴} “It’s a \RV{}{half}{-}\RVB{}{built}{}{build} \RV{}{house}{.} What did you \RV{}{expect}{?”} One \AG{}{pen}{,} one \AG{}{whole}{} night, all the \AG{}{holiday}{} \AG{}{homework}{} done. It \AGB{}{seems}{}{seem} \RV{}{impossible}{,} but it \RVB{}{happens}{}{happen} \RV{}{every}{} year.

\PARA{8}The last two days before school are also a \RV{}{real}{} \RV{}{test}{} of \RV{}{friendship}{.} \SMK{y}If your \AG{}{kid}{} is \RV{}{weak}{} at \RV{}{English}{} but \NS{}{strong}{}{新义}{}{/strɒŋ/}{adj.}{擅长的} at \RV{}{maths}{,} and has a friend who is \PH{the\ other\ way\ \RV{}{round}{,}}{the other way round}{正好相反} the two will \RV{}{probably}{} work \RV{}{together}{.}\EMK{y}{3} 
\YI{3}{如果你家孩子英语差、数学好，又有个朋友正好相反，两人多半会合作。}
\PH{All\ in\ all,}{all in all}{总之} the \RV{}{goal}{} is: \RV{“}{Teacher}{,} you can say I didn’t \PH{\RV{}{try}{}\ my\ best,}{try one's best}{尽最大努力} but you can’t say I didn’t do my \AG{}{homework}{!”}

\PARA{9}\RV{}{Hard}{} times \RV{}{bring}{} new \RVB{}{ideas}{.}{idea} When \RVB{}{visiting}{}{visit} \RVB{}{relatives}{,}{relative} \AG{}{kids}{} \AG{}{may}{} \RV{}{borrow}{} an older \NW{}{cousin}{’s}{/ˈkʌzn/}{n.[C]}{表哥或表姐} \AGB{}{finished}{}{finish} \AG{}{homework}{} from years before. They \AG{}{may}{} \RV{}{cry}{} \RV{}{until}{} their \AG{}{parents}{} help. I did that \RV{}{myself}{} \PH{\RV{}{as}{}\ a\ \AG{}{kid}{:}}{as a kid}{小时候} at \NW{}{midnight}{}{/ˈmɪdnaɪt/}{n.[U]}{午夜} I \AG{}{still}{} \AG{}{could}{n’t} \AG{}{finish}{,} so I \AGB{}{cried}{}{cry} \AG{}{until}{} my mum helped me write it. And \PH{if\ all\ \RV{}{else}{}\ \RVB{}{fails}{,}{fail}}{if all else fails}{万不得已时} the \AG{}{homework}{} \RV{}{simply}{} goes \RV{}{missing}{.} It got \RV{}{wet}{,} the \NW{}{dog}{}{/dɒɡ/}{n.[C]}{狗} ate it, the \AG{}{house}{} \PH{\RVB{}{caught}{}{catch}\ \RV{}{fire}{,}}{catch fire}{着火} or it was \RVB{}{left}{}{leave} on the \AG{}{underground}{.}

\PARA{10}For \AG{}{parents}{,} the night before school is \RV{}{full}{} of \RV{}{pressure}{.} The \AG{}{kid}{} is \RV{}{slow}{,} can’t do the \RVB{}{questions}{}{question} and can’t \PH{\RV{}{keep}{}\ his\ \RV{}{mind}{}\ on}{keep one's mind on}{专心于} the work. \AG{}{Parents}{} \PH{are\ \NW{}{afraid}{}{/əˈfreɪd/}{adj.}{害怕的}\ of}{be afraid of}{害怕……} a \RV{}{call}{} from the \PH{\NS{}{class}{}{新义}{}{/klɑːs/}{n.[C]}{班级}\ \AG{}{teacher}{.}}{class teacher}{班主任} \AG{}{Late}{} at night, many of us ask \NW{}{ourselves}{,}{/ˌaʊəˈselvz/}{pron.}{我们自己} “What did I do \NS{}{wrong}{?}{新词性}{}{/rɒŋ/}{adv.}{错；不对} Why are we doing this again?”

\PARA{11}Why does this \RV{}{keep}{} \AGB{}{happening}{?}{happen} A \RV{}{long}{} \AG{}{holiday}{} makes \AG{}{kids}{} too \PH{\RV{}{sure}{}\ of\ \RV{}{themselves}{:}}{be sure of oneself}{自信} “So much time! A \NS{}{bright}{}{新义}{}{/braɪt/}{adj.}{聪明的} \AG{}{kid}{} like me can \AG{}{finish}{} it \RV{}{easily}{.”} This \RV{}{kind}{} of \RV{}{thinking}{} \RVB{}{shows}{}{show} they have no \RV{}{idea}{} how \RV{}{fast}{} they \RV{}{really}{} work. \AG{}{Parents}{} can help them \PH{\NS{}{note}{}{新词性}{}{/nəʊt/}{vt.}{记录}\ down}{note down}{记下} how \AG{}{long}{} things take: two \NWB{}{pages}{}{page}{/peɪdʒ/}{n.[C]}{页} of new \NS{}{characters}{,}{新义}{character}{/ˈkærəktə(r)/}{n.[C]}{汉字；字} ten \AG{}{maths}{} \AG{}{questions}{,} or \PH{one\ \NS{}{piece}{}{新义}{}{/piːs/}{n.[C]}{一篇}\ of}{a piece of}{一篇（件）} \NS{}{writing}{.}{新义}{}{/ˈraɪtɪŋ/}{n.[U]}{作文；文章}

\PARA{12}\SMK{y}The \AGB{}{longer}{}{long} the \AG{}{holiday}{} goes on, the more \RV{}{everyone}{} plays.\EMK{y}{4} 
\YI{4}{假期过得越久，大家就玩得越多。}
\AG{}{Even}{} the \PH{\RVB{}{grown}{-}{grow}ups}{grown-up}{成年人} only want to eat, \RV{}{drink}{} and have \RV{}{fun}{.} Then \RV{}{try}{} to get the \AG{}{kids}{} to sit down and do their \AG{}{homework}{.} \SMK{g}What a \RV{}{challenge}{} that is!\EMK{g}{1} 
\GRAM{1}{\GITEM{感叹句}{感叹句用来表达强烈的感情，常用 what 或 how 引导。}{What a challenge 放句首，后接 that is，比 That is a challenge 语气强得多。}}{}
Some \AG{}{holiday}{} \AG{}{homework}{} is \RV{}{quite}{} \RV{}{difficult}{,} but \RV{}{with}{} so much time, \AG{}{kids}{} don’t \RV{}{feel}{} any \NS{}{hurry}{.}{新词性}{}{/ˈhʌri/}{n.[U]}{匆忙} A friend of \RV{}{mine}{} \RV{}{often}{} sees her \RV{}{daughter}{} in \RV{}{Grade}{} Two \PH{\NWB{}{staring}{}{stare}{/steə(r)/}{vi.}{盯着看}\ at}{stare at}{盯着看} a \AG{}{hard}{} \AG{}{question}{} and \NS{}{dreaming}{.}{新义}{dream}{/driːm/}{vi.}{走神；发呆} She wants to \RV{}{shout}{,} \AG{“}{Wake}{} up! \AG{}{Hard}{} \AG{}{questions}{} don’t \RV{}{disappear}{} \NS{}{like}{}{新义}{}{/laɪk/}{conj.}{像……那样} \AGB{}{holidays}{}{holiday} do!” So do the \AG{}{difficult}{} work \RV{}{early}{.}

\PARA{13}\RV{}{Sometimes}{,} the more \AG{}{parents}{} help, the \AGB{}{slower}{}{slow} \AG{}{kids}{} get. “If I don’t \RV{}{hear}{} the \AG{}{homework}{} \RV{}{clearly}{,} Mum will ask in the \AG{}{parents}{’} \PN{}{WeChat}{}{/ˈwiːtʃæt/}{n.}{微信} \RV{}{group}{.} If I don’t do it today, Mum will \RV{}{push}{} me \RV{}{tomorrow}{.”} When \AG{}{parents}{} \NS{}{catch}{}{新义}{}{/kætʃ/}{vt.}{接住；包办} \RV{}{everything}{,} \AG{}{kids}{} \RV{}{never}{} \RV{}{learn}{} to take \RV{}{care}{} of their \AG{}{own}{} work.

\PARA{14}Some \AG{}{kids}{} sit at the \AG{}{desk}{} for two \AG{}{hours}{} but work for only \NW{}{fifteen}{}{/ˌfɪfˈtiːn/}{num.}{十五} \RVB{}{minutes}{.}{minute} I \RV{}{would}{} \NW{}{advise}{}{/ədˈvaɪz/}{vt.}{建议} the \RV{}{usual}{} things: a \RV{}{quiet}{,} \RV{}{tidy}{} \RV{}{place}{} to \RV{}{study}{;} a \RV{“}{focus}{”} \PH{\NS{}{alarm}{}{新义}{}{/əˈlɑːm/}{n.[C]}{警报器；闹铃}\ \RV{}{clock}{,}}{alarm clock}{闹钟} so \AG{}{kids}{} \RV{}{pay}{} \RV{}{attention}{} for a \RV{}{short}{} time \AG{}{until}{} it \PH{goes\ \RV{}{off}{;}}{go off}{（闹钟）响} and \RV{}{clear}{} \RVB{}{rules}{,}{rule} like no \NS{}{play}{}{新义}{}{/pleɪ/}{n.[U]}{玩耍} \AG{}{until}{} the work is done.

\PARA{15}But I think one more thing \RVB{}{matters}{:}{matter} helping \AG{}{kids}{} \AG{}{enjoy}{} the work and \RV{}{believe}{} in \RV{}{themselves}{.} \SMK{y}If \AG{}{homework}{} were \RV{}{as}{} \NW{}{attractive}{}{/əˈtræktɪv/}{adj.}{有吸引力的} \RV{}{as}{} \AG{}{play}{,} \AG{}{kids}{} \AG{}{would}{} not \RV{}{leave}{} it all to the last night.\EMK{y}{5} 
\YI{5}{要是作业像玩一样有吸引力，孩子们就不会全都拖到最后一晚了。}
At the start of the \AG{}{holiday}{,} \AG{}{parents}{} and \AG{}{kids}{} can put the \AG{}{homework}{} \PH{in\ \RV{}{order}{}\ of}{in order of}{按……的顺序} what the child \RV{}{most}{} wants to do. \RV{}{Once}{} \RV{}{part}{} of it is done, the \RV{}{rest}{} \RVB{}{feels}{}{feel} \RVB{}{lighter}{,}{light} and \PH{\RV{}{step}{}\ by\ \AG{}{step}{,}}{step by step}{逐步地} the child’s \RV{}{confidence}{} \RVB{}{grows}{.}{grow}

\PARA{16}In the \RV{}{end}{,} don’t \PH{be\ too\ \AG{}{hard}{}\ on}{be hard on}{对……严厉} your \AG{}{kids}{,} or on \RV{}{yourself}{.} And if the work \AG{}{still}{} isn’t \AGB{}{finished}{,}{finish} let the child \RV{}{face}{} the \RV{}{result}{} \RV{}{alone}{} \PH{this\ \RV{}{once}{.}}{this once}{就这一次} That’s not a bad \NS{}{lesson}{}{新义}{}{/ˈlesn/}{n.[C]}{教训} \RV{}{either}{.}
\end{lessonbody}
\sourceline{Adapted from: “开学前一晚，在赶寒假作业现场，我感受了一把什么是中国速度……”, 果壳 (Guokr), 2024-02-27.}
\SelfTest{69}{\TW{complain}
\TW{skiing}
\TW{desk}
\TW{skating}
\TW{skate}
\TW{carefully}
\TW{writing}
\TW{cousin}
\TW{midnight}
\TW{dog}
\TW{afraid}
\TW{ourselves}
\TW{page}
\TW{stare}
\TW{fifteen}
\TW{advise}
\TW{attractive}
\TW{report\,\nstag{新词性}}
\TW{rush\,\nstag{新义}}
\TW{race\,\nstag{新义}}
\TW{treat\,\nstag{新义}}
\TW{hard\,\nstag{新义}}
\TW{picture\,\nstag{新词性}}
\TW{underground\,\nstag{新词性}}
\TW{strong\,\nstag{新义}}
\TW{class\,\nstag{新义}}
\TW{wrong\,\nstag{新词性}}
\TW{bright\,\nstag{新义}}
\TW{note\,\nstag{新词性}}
\TW{character\,\nstag{新义}}
\TW{piece\,\nstag{新义}}
\TW{writing\,\nstag{新义}}
\TW{hurry\,\nstag{新词性}}
\TW{dream\,\nstag{新义}}
\TW{like\,\nstag{新义}}
\TW{catch\,\nstag{新义}}
\TW{alarm\,\nstag{新义}}
\TW{play\,\nstag{新义}}
\TW{lesson\,\nstag{新义}}
\TW{report live}
\TW{use up}
\TW{work wonders}
\TW{sit up}
\TW{half-finished}
\TW{fight for one's life}
\TW{race against time}
\TW{ice skating}
\TW{make time}
\TW{answer back}
\TW{the other way round}
\TW{all in all}
\TW{try one's best}
\TW{as a kid}
\TW{if all else fails}
\TW{catch fire}
\TW{keep one's mind on}
\TW{be afraid of}
\TW{class teacher}
\TW{be sure of oneself}
\TW{note down}
\TW{a piece of}
\TW{grown-up}
\TW{stare at}
\TW{alarm clock}
\TW{go off}
\TW{in order of}
\TW{step by step}
\TW{be hard on}
\TW{this once}}
\par\write-1{LESSONEND 30 \thepage}
